% =====================================================================
%  Chapter 6 — Size Conventions and the Albedo-Temperature Degeneracy
% =====================================================================

\chapter{Size Conventions and the Albedo-Temperature Degeneracy}
\label{ch:size-conventions}

Before the energy of an impact can be computed, the size of the object has to
be fixed, and fixing the size of a small body is harder than it looks. This
chapter works through the two available routes to a diameter --- the
reflected-light route and the thermal route --- and shows that each carries an
assumption that is not measured. The degeneracy is not a defect of the
instruments. It is a property of the information available.

\section{The absolute magnitude and what it means}

A body's absolute magnitude \(H\) is the magnitude it would have at one
astronomical unit from the Sun and one astronomical unit from the observer,
at zero phase angle. It is a measure of the reflected light the body would
return if it were seen under standard conditions, and it is the only
photometric quantity the catalogue reliably carries.

The relation between \(H\), the geometric albedo \(p\) and the effective
diameter \(D\) is
\begin{equation}
D = \frac{1329}{\sqrt{p}}\,10^{-H/5}\ \ \mathrm{km},
\label{eq:hdp2}
\end{equation}
which follows from the definition of the albedo and the geometry of the
diffuse-reflection approximation. For a body of a given albedo, the diameter
is set by the brightness; for a body of a given brightness, the diameter is
set by the albedo.

The albedo is a composition measurement. A dark carbonaceous surface reflects
a few per cent; a bright silicate surface reflects several tens of per cent.
Changing the assumed albedo from 0.10 to 0.20 changes the diameter by a
factor of \(1/\sqrt{2}\), which is a factor of \(2.8\) in volume and a factor
of \(2.8\) in mass.

That is already a large effect, and it is only the first of two.

\section{The thermal route, and the beaming parameter}

The second route measures the body's own heat. A small body absorbs sunlight
and reradiates it, and the flux it emits in the thermal infrared is a
function of its temperature, its emissivity and its projected area --- which
is a function of its diameter.

The equilibrium temperature\idx{Equilibrium temperature} follows from a balance. The absorbed flux at
heliocentric distance \(r\), for a body with geometric albedo \(A\) and
emissivity \(\varepsilon\), is \((1-A)S_0/r^{2}\) per unit area, where \(S_0\)
is the solar constant at one astronomical unit. The emitted flux is
\(\varepsilon\sigma T^{4}\). Balancing them and inserting the solar constant
gives the familiar result
\begin{equation}
T_{\mathrm{eq}} = 278\ \mathrm{K}\ \left(\frac{1\ \mathrm{AU}}{r}\right)^{1/2}
\label{eq:teq}
\end{equation}
for \(A = 0.1\) and \(\varepsilon = 0.9\), and a factor \(((1-A)/\varepsilon)^{1/4}\)
for other values. The temperature does not depend on the size, which is the
reason the technique works.

The observed flux, however, does. The total thermal flux from a body of
diameter \(D\) is proportional to \(D^{2}\), and the relationship between the
observed flux and the body's own emission depends on how the heat is
distributed over the surface. That distribution is summarised by the
\emph{beaming parameter} \(\eta\): it is the ratio of the total emitted power
to the power that a smooth sphere would emit at the same subsolar
temperature. For a body whose surface radiates efficiently the beaming
parameter approaches unity; for a rough surface with self-shadowing it is
smaller. The physically allowed range is \(0 < \eta \leq 1\).

The observed flux therefore scales as
\begin{equation}
F_{\mathrm{obs}} \propto \eta\,D^{2},
\label{eq:eta}
\end{equation}
and \eqref{eq:eta} is the degeneracy. The diameter and the beaming parameter
cannot be separated by a single thermal band: the measurement determines the
product, and the diameter is fixed only by assuming a beaming parameter.

\section{What the beaming parameter costs, and how it is detected}

The assumption is testable, because a beaming parameter above unity is
unphysical, and a fit that returns one is telling you that the model is
wrong rather than that the object is strange. The 2025 JWST analysis of
2024~YR4 is an exact illustration.

\begin{table}[!htbp]
\centering
\small
\caption{The three thermal models fitted to 2024 YR4, and what the beaming
parameter did to each.}
\label{tab:thermal-models}
\begin{tabularx}{\textwidth}{@{}P{1.14in}P{0.74in}P{0.82in}Y@{}}
\toprule
\rowcolor{panel}
\textbf{Model} & \textbf{Diameter} & \textbf{Albedo} & \textbf{Status} \\
\midrule
NEATM, standard & 60.6 m & --- & requires a beaming parameter of 3.06,
which is unphysical; rejected \\
\rowcolor{panel}
Fast rotating model & 51.1 m & 0.16 & a coherent rotation model, credible \\
Thermophysical model & 60.8 m & --- & the published best fit; a
non-rotating body with a thermal inertia\idx{Thermal inertia} \\
\bottomrule
\end{tabularx}
\end{table}

The two credible models differ by 19 per cent in diameter: 51.1 against
60.8 metres. The photometry\idx{Photometry} that produced them is excellent. The difference
between them is entirely a statement about how the heat is distributed, and
therefore about how the surface is shaped.

\begin{deeper}{Why the thermophysical model is not simply better}
The thermophysical model treats the body as having a finite thermal inertia\idx{Thermal inertia},
so the surface temperature varies through the rotation and the night is
cooler than the day. That is physically more realistic, and it is the source
of the high thermal inertia\idx{Thermal inertia} and low surface roughness reported for
2024~YR4. But realism in the forward model is not the same as freedom from
assumption: the thermophysical model still has to be supplied with a spin
state, a thermal inertia and a roughness distribution, and it delivers a
diameter that is 19 per cent larger than the one the coherent-rotation model
gives. The correct reading is not that one model is right. It is that the
model choice is worth more than the measurement error, and that a
measurement which does not depend on the model --- a resolved image --- is
worth a great deal.
\end{deeper}

\section{The photometric route, and its own assumption}

The third route is the one the field uses most, because it is the cheapest.
A lightcurve gives the object's rotation period and the projected extent of
its brightness variation; combined with the absolute magnitude, an inversion
produces a shape model, and from the model a volume, a diameter and an
albedo.

The assumption here is the convexity. A lightcurve constrains the body's
projected outline in two dimensions as it rotates, and the standard inversion
returns the convex hull\idx{Convex hull} of the constraint set --- a convex body, by
construction, whatever the body actually is. This is the subject of
Chapter 17.

\section{Why the degeneracy is the practical problem}

Taken together, the three routes have this structure. The photometric route
assumes convexity and returns an albedo as an output rather than an input.
The thermal route assumes a beaming parameter and cannot detect it from a
single band. The absolute magnitude route assumes an albedo and has no way
to check it.

\begin{table}[!htbp]
\centering
\small
\caption{The three routes to a diameter, the assumption each one makes, and
the size of the resulting uncertainty in the energy.}
\label{tab:routes}
\begin{tabularx}{\textwidth}{@{}P{1.36in}P{1.68in}P{1.23in}@{}}
\toprule
\rowcolor{panel}
\textbf{Route} & \textbf{Assumption} & \textbf{Energy factor} \\
\midrule
Thermal, single model & beaming parameter & 1.0, by definition \\
\rowcolor{panel}
Thermal, model spread & choice of model & 1.68 for 2024 YR4 \\
Absolute magnitude & albedo & up to 2.8 \\
\rowcolor{panel}
Photometric inversion & convexity & of order 1.1 to 1.3 \\
Resolved imaging & none of the above & bounded by the imaging \\
\bottomrule
\end{tabularx}
\end{table}

The last row is the reason this book exists. A resolved image does not need
an albedo, a beaming parameter or a convexity assumption. It reports the
projected area at the epoch, and the projected area is a measurement. What it
does not directly report is the volume, because a projection of a rotating
body is not a projection of its shape unless the body is convex --- but it
reports the two-dimensional geometry that the volume has to be built on, and
it reports it without a model.

\begin{keytakeaways}
\item A diameter cannot be obtained from photometry\idx{Photometry} without an assumption
about the albedo, and cannot be obtained from a single thermal band without
an assumption about the beaming parameter.
\item The beaming parameter is bounded by unity, so a fit returning 3.06 is a
detectable model failure rather than a measurement.
\item The two credible thermal models for 2024 YR4 differ by 19 per cent in
diameter, which is a factor of 1.68 in energy at fixed density.
\item The measured values in \tabref{tab:routes} are the model choices, and
the resolved imaging row is the only one with no model in it.
\end{keytakeaways}
